\RequirePackage[l2tabu]{nag}		% Warns for incorrect (obsolete) LaTeX usage
%
%
% File: memoirthesis.tex
% Author: Abhimanyu Bambhaniya 
% Author: Nikhil Kumar 
% Description: Contains the thesis template using memoir class,
% which is mainly based on book class but permits better control of
% chapter styles for example. This template is an adaptation and
% modification of Oscar's.
%
% Memoir is a flexible class for typesetting poetry, fiction,
% non-fiction and mathematical works as books, reports, articles or
% manuscripts. CTAN repository is found at:
% http://www.ctan.org/tex-archive/macros/latex/contrib/memoir/
%
%
% UoB guidelines for thesis presentation were found at:
% http://www.bris.ac.uk/esu/pg/pgrcop11-12topic.pdf#page=49
%
% Memoir class loads useful packages by default (see manual).
\documentclass[oneside,a4paper,12pt,reqno,openbib]{memoir} %add 'draft' to turn draft option on (see below)

% Adding metadata:
\usepackage{listings}
\usepackage{relsize}
\usepackage[table]{xcolor}
\usepackage{subcaption}
\usepackage{algorithm,algorithmicx,algpseudocode}
%\renewcommand{\algorithmcfname}{ALGORITHM}
%\SetAlFnt{\small}
%\SetAlCapFnt{\small}
%\SetAlCapNameFnt{\small}
%\SetAlCapHSkip{0pt}
%\IncMargin{-\parindent}

\setlength{\parskip}{1em}

\usepackage{datetime}
\usepackage{ifpdf}
\ifpdf
\pdfinfo{
   /Author (Abhimanyu Bambhaniya)
   /Title (Undergraduate Dissertation)
   /CreationDate (D:\pdfdate)
}
\fi
% When draft option is on.
\ifdraftdoc
	\usepackage{draftwatermark}				%Sets watermarks up.
	\SetWatermarkScale{0.3}
	\SetWatermarkText{\bf Draft: \today}
\fi
%
% Declare figure/table as a subfloat.
\newsubfloat{figure}
\newsubfloat{table}
% Better page layout for A4 paper, see memoir manual.
\settrimmedsize{297mm}{210mm}{*}
\setlength{\trimtop}{0pt}
\setlength{\trimedge}{\stockwidth}
\addtolength{\trimedge}{-\paperwidth}
\settypeblocksize{634pt}{448.13pt}{*}
\setulmargins{4cm}{*}{*}
\setlrmargins{*}{*}{0.6}
%\setlrmarginsandblock{1.5in}{1in}{*}
\setmarginnotes{17pt}{51pt}{\onelineskip}
\setheadfoot{\onelineskip}{2\onelineskip}
\setheaderspaces{*}{2\onelineskip}{*}
\checkandfixthelayout
%
\frenchspacing
% Font with math support: New Century Schoolbook
\usepackage{fouriernc}
\usepackage[T1]{fontenc}
%
% UoB guidelines:
%
% Text should be in double or 1.5 line spacing, and font size should be
% chosen to ensure clarity and legibility for the main text and for any
% quotations and footnotes. Margins should allow for eventual hard binding.
%
% Note: This is automatically set by memoir class. Nevertheless \OnehalfSpacing
% enables double spacing but leaves single spaced for captions for instance.
\OnehalfSpacing
%
% Sets numbering division level
\setsecnumdepth{subsection}
\maxsecnumdepth{subsubsection}
%
% Chapter style (taken and slightly modified from Lars Madsen Memoir Chapter
% Styles document
\usepackage{calc,soul,fourier}
\makeatletter
\newlength\dlf@normtxtw
\setlength\dlf@normtxtw{\textwidth}
\newsavebox{\feline@chapter}
\newcommand\feline@chapter@marker[1][4cm]{%
	\sbox\feline@chapter{%
		\resizebox{!}{#1}{\fboxsep=1pt%
			\colorbox{gray}{\color{white}\thechapter}%
		}}%
		\rotatebox{90}{%
			\resizebox{%
				\heightof{\usebox{\feline@chapter}}+\depthof{\usebox{\feline@chapter}}}%
			{!}{\scshape\so\@chapapp}}\quad%
		\raisebox{\depthof{\usebox{\feline@chapter}}}{\usebox{\feline@chapter}}%
}
\newcommand\feline@chm[1][4cm]{%
	\sbox\feline@chapter{\feline@chapter@marker[#1]}%
	\makebox[0pt][c]{% aka \rlap
		\makebox[1cm][r]{\usebox\feline@chapter}%
	}}
\makechapterstyle{daleifmodif}{
	\renewcommand\chapnamefont{\normalfont\Large\scshape\raggedleft\so}
	\renewcommand\chaptitlefont{\normalfont\Large\bfseries\scshape}
	\renewcommand\chapternamenum{} \renewcommand\printchaptername{}
	\renewcommand\printchapternum{\null\hfill\feline@chm[2.5cm]\par}
	\renewcommand\afterchapternum{\par\vskip\midchapskip}
	\renewcommand\printchaptertitle[1]{\color{black}\chaptitlefont\raggedleft ##1\par}
}
\makeatother
\chapterstyle{daleifmodif}
%
% UoB guidelines:
%
% The pages should be numbered consecutively at the bottom centre of the
% page.
\makepagestyle{myvf}
\makeoddfoot{myvf}{}{\thepage}{}
\makeevenfoot{myvf}{}{\thepage}{}
\makeheadrule{myvf}{\textwidth}{\normalrulethickness}
\makeevenhead{myvf}{\small\textsc{\leftmark}}{}{}
\makeoddhead{myvf}{}{}{\small\textsc{\rightmark}}
\pagestyle{myvf}
%
% Oscar's command (it works):
% Fills blank pages until next odd-numbered page. Used to emulate single-sided
% frontmatter. This will work for title, abstract and declaration. Though the
% contents sections will each start on an odd-numbered page they will
% spill over onto the even-numbered pages if extending beyond one page
% (hopefully, this is ok).
\newcommand{\clearemptydoublepage}{\newpage{\thispagestyle{empty}\cleardoublepage}}
%
%
% Creates indexes for Table of Contents, List of Figures, List of Tables and Index
\makeindex
% \printglossaries below creates a list of abbreviations. \gls and related
% commands are then used throughout the text, so that latex can automatically
% keep track of which abbreviations have already been defined in the text.
%
% The import command enables each chapter tex file to use relative paths when
% accessing supplementary files. For example, to include
% chapters/brewing/images/figure1.png from chapters/brewing/brewing.tex we can
% use
% \includegraphics{images/figure1}
% instead of
% \includegraphics{chapters/brewing/images/figure1}
\usepackage{import}

% Add other packages needed for chapters here. For example:
\usepackage{lipsum}					%Needed to create dummy text
\usepackage{amsfonts} 					%Calls Amer. Math. Soc. (AMS) fonts
\usepackage[centertags]{amsmath}			%Writes maths centred down
\usepackage{stmaryrd}					%New AMS symbols
\usepackage{amssymb}					%Calls AMS symbols
\usepackage{amsthm}					%Calls AMS theorem environment
\usepackage{newlfont}					%Helpful package for fonts and symbols
\usepackage{mathtools}
\DeclarePairedDelimiter\abs{\lvert}{\rvert}%
\usepackage{layouts}					%Layout diagrams
\usepackage{graphicx}					%Calls figure environment
\graphicspath{{images/}}
\usepackage{wrapfig}
\usepackage{pdfpages}
\usepackage{lscape}
\usepackage{pdflscape}
\usepackage{rotating}
\usepackage{epstopdf}
%\usepackage{subfigure}
\usepackage{longtable,rotating}			%Long tab environments including rotation.
\usepackage[applemac]{inputenc}			%Needed to encode non-english characters
\usepackage{gensymb}					
\usepackage{colortbl}					%Makes coloured tables
\usepackage{wasysym}					%More math symbols
\usepackage{mathrsfs}					%Even more math symbols
\usepackage{float}						%Helps to place figures, tables, etc.
\usepackage{verbatim}					%Permits pre-formated text insertion
\usepackage{upgreek }					%Calls other kind of greek alphabet
\usepackage{latexsym}					%Extra symbols
\usepackage[square,numbers,
		     sort&compress]{natbib}		%Calls bibliography commands
\usepackage{url}						%Supports url commands
\usepackage{etex}						%eTeXÕs extended support for counters
\usepackage{fixltx2e}					%Eliminates some in felicities of the
									%original LaTeX kernel
\usepackage[spanish,english]{babel}		%For languages characters and hyphenation
\usepackage{color}                    				%Creates coloured text and background
\usepackage[colorlinks=true,
		     allcolors=black]{hyperref}              %Creates hyperlinks in cross references
\usepackage{memhfixc}					%Must be used on memoir document
									%class after hyperref
\usepackage{enumerate}					%For enumeration counter
\usepackage{footnote}					%For footnotes
\usepackage{microtype}					%Makes pdf look better.
\usepackage{rotfloat}					%For rotating and float environments as tables,
									%figures, etc.
\usepackage{alltt}						%LaTeX commands are not disabled in
									%verbatim-like environment
\usepackage[version=0.96]{pgf}			%PGF/TikZ is a tandem of languages for producing vector graphics from a
\usepackage{tikz}	
\usetikzlibrary{calc}
\usetikzlibrary{arrows,shapes,snakes,
		       automata,backgrounds,
		       petri,topaths}				%To use diverse features from tikz		
%							
%Reduce widows  (the last line of a paragraph at the start of a page) and orphans
% (the first line of paragraph at the end of a page)
\widowpenalty=1000
\clubpenalty=1000
%
% New command definitions for my thesis
%
\newcommand{\keywords}[1]{\par\noindent{\small{\bf Keywords:} #1}} %Defines keywords small section
\newcommand{\parcial}[2]{\frac{\partial#1}{\partial#2}}                             %Defines a partial operator
\newcommand{\vectorr}[1]{\mathbf{#1}}                                                        %Defines a bold vector
\newcommand{\vecol}[2]{\left(                                                                         %Defines a column vector
	\begin{array}{c}
		\displaystyle#1 \\
		\displaystyle#2
	\end{array}\right)}
\newcommand{\mados}[4]{\left(                                                                       %Defines a 2x2 matrix
	\begin{array}{cc}
		\displaystyle#1 &\displaystyle #2 \\
		\displaystyle#3 & \displaystyle#4
	\end{array}\right)}
\newcommand{\pgftextcircled}[1]{                                                                    %Defines encircled text
    \setbox0=\hbox{#1}%
    \dimen0\wd0%
    \divide\dimen0 by 2%
    \begin{tikzpicture}[baseline=(a.base)]%
        \useasboundingbox (-\the\dimen0,0pt) rectangle (\the\dimen0,1pt);
        \node[circle,draw,outer sep=0pt,inner sep=0.1ex] (a) {#1};
    \end{tikzpicture}
}
\newcommand{\range}[1]{\textnormal{range }#1}                                             %Defines range operator
\newcommand{\innerp}[2]{\left\langle#1,#2\right\rangle}                                 %Defines inner product
\newcommand{\prom}[1]{\left\langle#1\right\rangle}                                         %Defines average operator
\newcommand{\tra}[1]{\textnormal{tra} \: #1}                                                       %Defines trace operator
\newcommand{\sign}[1]{\textnormal{sign\,}#1}                                                   %Defines sign operator
\newcommand{\sech}[1]{\textnormal{sech} #1}                                                  %Defines sech
\newcommand{\diag}[1]{\textnormal{diag} #1}                                                    %Defines diag operator
\newcommand{\arcsech}[1]{\textnormal{arcsech} #1}                                       %Defines arcsech
\newcommand{\arctanh}[1]{\textnormal{arctanh} #1}                                         %Defines arctanh
%Change tombstone symbol
\newcommand{\blackged}{\hfill$\blacksquare$}
\newcommand{\whiteged}{\hfill$\square$}
\newcounter{proofcount}
\renewenvironment{proof}[1][\proofname.]{\par
 \ifnum \theproofcount>0 \pushQED{\whiteged} \else \pushQED{\blackged} \fi%
 \refstepcounter{proofcount}
 \normalfont
 \trivlist
 \item[\hskip\labelsep
       \itshape
   {\bf\em #1}]\ignorespaces
}{%
 \addtocounter{proofcount}{-1}
 \popQED\endtrivlist
}
%
%
% New definition of square root:
% it renames \sqrt as \oldsqrt
\let\oldsqrt\sqrt
% it defines the new \sqrt in terms of the old one
\def\sqrt{\mathpalette\DHLhksqrt}
\def\DHLhksqrt#1#2{%
\setbox0=\hbox{$#1\oldsqrt{#2\,}$}\dimen0=\ht0
\advance\dimen0-0.2\ht0
\setbox2=\hbox{\vrule height\ht0 depth -\dimen0}%
{\box0\lower0.4pt\box2}}
%
% My caption style
\newcommand{\mycaption}[2][\@empty]{
	\captionnamefont{\scshape}
	\changecaptionwidth
	\captionwidth{0.9\linewidth}
	\captiondelim{.\:}
	\indentcaption{0.75cm}
	\captionstyle[\centering]{}
	\setlength{\belowcaptionskip}{10pt}
	\ifx \@empty#1 \caption{#2}\else \caption[#1]{#2}
}

\newcommand*\vtick{\textsc{\char13}}
\newcommand*\expandedname{Twitter News Feed}
\newcommand*\shortname{TNF}

%
% My subcaption style
\newcommand{\mysubcaption}[2][\@empty]{
	\subcaptionsize{\small}
	\hangsubcaption
	\subcaptionlabelfont{\rmfamily}
	\sidecapstyle{\raggedright}
	\setlength{\belowcaptionskip}{10pt}
	\ifx \@empty#1 \subcaption{#2}\else \subcaption[#1]{#2}
}
%
%An initial of the very first character of the content
\usepackage{lettrine}
\newcommand{\initial}[1]{%
	\lettrine[lines=3,lhang=0.33,nindent=0em]{
		\color{gray}
     		{\textsc{#1}}}{}}
%
% Theorem styles used in my thesis
%
\theoremstyle{plain}
\newtheorem{theo}{Theorem}[chapter]
\theoremstyle{plain}
\newtheorem{prop}{Proposition}[chapter]
\theoremstyle{plain}
\theoremstyle{definition}
\newtheorem{dfn}{Definition}[chapter]
\theoremstyle{plain}
\newtheorem{lema}{Lemma}[chapter]
\theoremstyle{plain}
\newtheorem{cor}{Corollary}[chapter]
\theoremstyle{plain}
\newtheorem{resu}{Result}[chapter]
%
% Hyphenation for some words
%
\hyphenation{res-pec-tively}
\hyphenation{mono-ti-ca-lly}
\hyphenation{hypo-the-sis}
\hyphenation{para-me-ters}
\hyphenation{sol-va-bi-li-ty}
%
%
\begin{document}
% UoB guidlines:
%
% Preliminary pages
%
% The five preliminary pages must be the Title Page, Abstract, Dedication
% and Acknowledgements, Author's Declaration and Table of Contents.
% These should be single-sided.
%
% Table of contents, list of tables and illustrative material
%
% The table of contents must list, with page numbers, all chapters,
 % sections and subsections, the list of references, bibliography, list of
% abbreviations and appendices. The list of tables and illustrations
% should follow the table of contents, listing with page numbers the
% tables, photographs, diagrams, etc., in the order in which they appear
% in the text.
%
\frontmatter
\pagenumbering{roman}
%
\begin{titlingpage}
\begin{SingleSpace}
\calccentering{\unitlength}
\begin{adjustwidth*}{\unitlength}{-\unitlength}
\vspace*{-1.5 cm}
\begin{center}
{\LARGE \textbf{Design of an Architecture for\\ In Memory Computation using 8T SRAM}}
\vspace{1mm}
{\Large \textbf{}}\\[4mm]
\vspace{0.5 cm}
{\large \emph{A Thesis\\ submitted towards the partial fulfillment of the\\ \vspace{1 mm} requirements for the award of the degree \\ \vspace{1 mm} of}}\\
\vspace{5 mm}
{\large \textbf{BACHELOR OF TECHNOLOGY} \vspace{3 mm} \\\emph{in}\\\vspace{3 mm} \textbf{ELECTRONICS AND COMMUNICATION ENGINEERING}}\\
\vspace{4 mm}
%{\large \textbf{(With specialization in Information Technology)}} \\
\vspace{2 mm}
{\large \emph{Submitted by}}\\
\vspace{5mm}

\begin{minipage}[t]{0.8\textwidth}
\begin{center} \large
{\large\textbf{Abhimanyu Bambhaniya (15116003)}}\\
{\large\textbf{Nikhil Kumar (15116036)}}\\
\end{center}
\end{minipage}

%---------------------------------------
\vspace{0 mm}
{\large \textit{Under the guidance of:}}\\
\vspace{3 mm}
{\large\textbf{Dr. Sudeb Dasgupta\\ \& \\ Dr. Anand Bulusu}}\\%---------------------------------------
\vspace{8 mm}
\includegraphics[height = 4cm, keepaspectratio]{IITR_Logo.jpg}\\
\vspace{8 mm}
{\large \textbf{Department of Electronics and Communication Engineering}\\}
\vspace{1mm}
{\large \textbf{Indian Institute of Technology Roorkee}}\\
\vspace{1mm}
{\large \textbf{Roorkee $-$ 247667 (India)}}\\
\vspace{1mm}
{\large \textbf{May, 2019}}
%\vspace{5 mm}
\end{center}
%\begin{flushright}
%{\small Word count: ten thousand and four}
%\end{flushright}
\end{adjustwidth*}
\end{SingleSpace}
\end{titlingpage} 
\clearemptydoublepage
%
%
% File: declaration.tex
% Author: V?ctor Bre?a-Medina
% Description: Contains the declaration page
%
% UoB guidelines:
%
% Author's declaration
%
% I declare that the work in this dissertation was carried out in accordance
% with the requirements of the University's Regulations and Code of Practice
% for Research Degree Programmes and that it has not been submitted for any
% other academic award. Except where indicated by specific reference in the
% text, the work is the candidate's own work. Work done in collaboration with,
% or with the assistance of, others, is indicated as such. Any views expressed
% in the dissertation are those of the author.
%
% SIGNED: .............................................................
% DATE:..........................
%
%\chapter*{Candidate's Declaration}
\noindent \hfill \textsc{\Large Student's Declaration}\\
\noindent\makebox[\linewidth]{\rule{\textwidth}{1pt}} 
\begin{SingleSpacing}
%\begin{quote}
\vspace{-5mm}
I declare that the work presented in this report with title ``\textbf{Design of an Architecture for In Memory Computation using 8T SRAM}" towards the fulfillment of the requirement for the award of the degree of \textbf{Bachelor of Technology} in \textbf{Electronics \& Communication Engineering} submitted in the \textbf{Department of Electronics \& Communication Engineering}, \textbf{Indian Institute of Technology, Roorkee}, India is an authentic record of my own work carried out during the period \textbf{from August 2018 to May 2019} under the supervision of \textbf{Dr. Sudeb Dasgupta}, Professor, Dept. of ECE, IIT Roorkee and \textbf{Dr. Anand Bulusu}, Associate Professor, Dept. of ECE, IIT Roorkee.
\noindent The content of this report has not been submitted by me for the award of any other degree of this or any other institute.

\vspace{3cm}
\noindent
\textsc{DATE : ......................... ~~~~~~~~~~~~~~~~~~~~~~~~~~~~~~~~~~SIGNED: ........................................}\\
\textsc{PLACE: .........................}
\nobreak
\hfill \textsc{(Abhimanyu Bambhaniya)} \\
\vspace{-4 mm}
\nobreak
\hfill \textsc{Enrollment No.: 15116003}
\\\\\\\\
\vspace{-4 mm}
\nobreak
\hfill\textsc{SIGNED: ........................................}\\\\
\vspace{-4 mm}
\nobreak
\hfill \textsc{(Nikhil Kumar)} \\\\
\vspace{-4 mm}
\nobreak
\hfill \textsc{Enrollment No.: 15116036}

\end{SingleSpacing}
\clearpage 
\clearemptydoublepage


%
% File: declaration.tex
% Author: V?ctor Bre?a-Medina
% Description: Contains the declaration page
%
% UoB guidelines:
%
% Author's declaration
%
% I declare that the work in this dissertation was carried out in accordance
% with the requirements of the University's Regulations and Code of Practice
% for Research Degree Programmes and that it has not been submitted for any
% other academic award. Except where indicated by specific reference in the
% text, the work is the candidate's own work. Work done in collaboration with,
% or with the assistance of, others, is indicated as such. Any views expressed
% in the dissertation are those of the author.
%
% SIGNED: .............................................................
% DATE:..........................
%
\chapter*{Supervisor's Certificate}
\begin{OnehalfSpacing}
%\begin{quote}
This is to certify that the statement made by the candidate is correct to the best of my knowledge and belief.

\vspace{2.25cm}


\noindent
\textsc{DATE : .........................~~~~~~~~~~~~~~~~~~~~~~~~~~~~~~~~~~~SIGNED: ........................................}

\nobreak
\hfill \textsc{(Dr. SUDEB DASGUPTA)}

\vspace{-4 mm}
\nobreak
\hfill \textsc{Professor}

\vspace{-4 mm}
\nobreak
\hfill \textsc{Dept. of ECE, IIT Roorkee}

\vspace{1.25cm}
\noindent
\textsc{DATE : .........................~~~~~~~~~~~~~~~~~~~~~~~~~~~~~~~~~~~SIGNED: ........................................}

\nobreak
\hfill \textsc{(Dr. ANAND BULUSU)}

\vspace{-4 mm}
\nobreak
\hfill \textsc{Associate Professor}

\vspace{-4 mm}
\nobreak
\hfill \textsc{Dept. of ECE, IIT Roorkee}

\end{OnehalfSpacing}
\clearpage 
\clearemptydoublepage
%

%
% file: dedication.tex
% author: V?ctor Bre?a-Medina
% description: Contains the text for thesis dedication
%

\chapter*{Acknowledgements}
\begin{OnehalfSpacing}
%\begin{center}
%\begin{minipage}{14.7 cm}
\initial I sincerely thank my B.Tech. Project advisors \textbf{Dr. Sudeb Dasgupta} and \textbf{Dr. Anand Bulusu} for their constant guidance and support throughout the project. 

I would like to acknowledge the contribution of \textbf{Mr. Dinesh Kushwaha}, research scholar in Microelectronics \& VLSI, Department of ECE, IIT Roorkee for his ideas and motivation. I am also thankful to all the research scholars of SMDP Lab, Department of ECE for helpful discussions at every stage of the project. I thank \textbf{Department of Electronics and Communication Engineering, IIT Roorkee} for providing the lab and other resources required for carrying out my project work. 

Overall, working on this project has been a knowledgeable experience and provided me with a great experience in the area of Digital VLSI Circuits and Architecture. 


%\end{minipage}
%\end{center}
\end{OnehalfSpacing}
\clearpage

\clearemptydoublepage
%
%
% File: abstract.tex
% Author: V?ctor Bre?a-Medina
% Description: Contains the text for thesis abstract
%
% UoB guidelines:
%
% Each copy must include an abstract or summary of the dissertation in not
% more than 300 words, on one side of A4, which should be single-spaced in a
% font size in the range 10 to 12. If the dissertation is in a language other
% than English, an abstract in that language and an abstract in English must
% be included.

\chapter*{Abstract}
%\begin{SingleSpace}
\begin{OnehalfSpacing}
\paragraph{}
\end{OnehalfSpacing}

For providing novel solutions to the problems of memory bottleneck in Von-Neumann Architecture, we are trying to modify the most basic memory element, the S-RAM cell and it’s peripherals, but without any substantial increase in area. As the industry of Machine
Learning, Artificial Intelligence and IOT are on the rise, these memories can help with In-Memory Computing operations thus reducing machine time.

We designed an Architecture using 4x8 8T SRAM array for In Memory Computation with a word length of 4 bits in 180nm technology (SCL PDK) at 1.8V Vdd. This architecture is capable of computing logical operations like NAND, NOR, XOR, NOT and ADD, SUB operation for two words without conventional gate designs. LEFT SHIFT and RIGHT SHIFT operations are also possible in our design that too without the use of Flip Flops. Unlike conventional memory decoders, the decoders for this circuit require two addresses for reading operands and one address for writing the result along with the Opcode.  

The simulations were done using Cadence and Python together, in which the python code converted the Assembly code into bit streams for the control signals in the architecture. So it can be used like a mini computer to perform small tasks, by just writing an Assembly code.  

All the design, simulations and tests were performed in Cadence Virtuoso ADE and Python.

\clearpage 
\clearemptydoublepage
%
\renewcommand{\contentsname}{Table of Contents}
\maxtocdepth{subsection}
\tableofcontents*
\addtocontents{toc}{\par\nobreak \mbox{}\hfill{\bf Page}\par\nobreak}
\clearemptydoublepage
%
%\listoftables
%\addtocontents{lot}{\par\nobreak\textbf{{\scshape Table} \hfill %Page}\par\nobreak}
%\clearemptydoublepage
%
\listoffigures
\addtocontents{lof}{\par\nobreak\textbf{{\scshape Figure} \hfill Page}\par\nobreak}
\clearemptydoublepage
%
% The bulk of the document is delegated to these chapter files in
% subdirectories.
\mainmatter
%
\import{chapters/chapter01/}{chap01.tex}
\clearemptydoublepage
%
\import{chapters/chapter02/}{chap02.tex}
\clearemptydoublepage
%
\import{chapters/chapter03/}{chap03.tex}
\clearemptydoublepage
%
\import{chapters/chapter04/}{chap04.tex}
\clearemptydoublepage
%
\import{chapters/chapter05/}{chap05.tex}
\clearemptydoublepage
%
\import{chapters/chapter06/}{chap06.tex}
\clearemptydoublepage
%
\import{chapters/chapter07/}{chap07.tex}
\clearemptydoublepage

\import{chapters/chapter08/}{chap08.tex}
\clearemptydoublepage
%

%
%
%
% Apparently the guidelines don't say anything about citations or
% bibliography styles so I guess we can use anything.
\backmatter
\bibliographystyle{ieeetr}
\refstepcounter{chapter}
\bibliography{thesisbiblio}
\clearemptydoublepage

\appendix
\def\code#1{\texttt{#1}}
\chapter{Appendix}
\label{chap:intro}
\vspace{-0.2 cm}



\section{ Python code for converting Assembly to Bit Stream}

\paragraph{}
\begin{lstlisting}

NO_OF_INSTRUCTIONS = 9
instruction_file = 	open("instruction.txt", "r") 
write_row_0 =	 	open("wr_add0.txt","w+")
write_row_1 =	 	open("wr_add1.txt","w+")
write_pre_0 =		open("wr_pre0.txt","w+")
write_pre_1 =		open("wr_pre1.txt","w+")
write_col_en0 =		open("write_col_en0","w+")
write_col_en1 =		open("write_col_en1","w+")
write_en=		open("write_en","w+")
Din_s =			open("Din_s","w+")
data1 =			open("data1.txt","w+")
data2 =			open("data2.txt","w+")
data3 =			open("data3.txt","w+")
data4 =			open("data4.txt","w+")

out =	 		open("out.txt","w+")
enW =	 		open("enW.txt","w+")
en1 =			open("en1.txt","w+")
en =			open("en.txt","w+")
readP =			open("readP.txt","w+")
enb =			open("enb.txt","w+")
writeP=			open("writeP.txt","w+")
RL =			open("RL.txt","w+")
readP2 =		open("readP2.txt","w+")
S =				open("S.txt","w+")
Sb =			open("Sb.txt","w+")


read_pre_0 =		open("read_pre0.txt","w+")
read_pre_1 =		open("read_pre1.txt","w+")
r_add0     =		open("r_add0.txt","w+")
r_add1     =		open("r_add1.txt","w+")
read_en	   = 		open("read_en","w+")
op_sel1	   =		open("op_sel1","w+")
op_sel2	   =		open("op_sel2","w+")
block_en1  =		open("block1_en","w+")
block_en2  =		open("block2_en","w+")
nor	   = 			open("nor","w+")

def write_defaults(bit):
	for i in range(bit):							#Row address
			write_row_1.write("0")
			write_row_0.write("0")
			write_pre_0.write("1")
			write_pre_1.write("1")	
			write_col_en0.write("0")
			write_col_en1.write("0")
			Din_s.write("0")					#Din mux select
			write_en.write("0")					#Row enable		
			data1.write("0")
			data2.write("0")
			data3.write("0")
			data4.write("0")	
			out.write("1")
			enW.write("0")
			en1.write("1")
			en.write("0")
			readP.write("0")
			enb.write("1")
			RL.write("1")
			writeP.write("0")
			readP2.write("0")
	return
def read_defaults(bit):
	for i in range(bit):							
			r_add0.write("0")
			r_add1.write("0")
			read_pre_0.write("1")
			read_pre_1.write("1")	
			op_sel1.write("0")
			op_sel2.write("0")
			nor.write("0")						
			read_en.write("0")							
			block_en1.write("0")
			block_en2.write("0")
				
	return
	
for index in range(NO_OF_INSTRUCTIONS):
	inst=instruction_file.next().replace(" ", "")
	print "This is inst %d" %(index)
	print(inst[0:4])
	print(inst[4])  
	print(inst[4:7]) 
	print(inst[7:10])
	print(inst[10:13])
	print(inst[13:17])
	if(inst[0:4]=="0000"):									#Takes 23 cycles
		bits=23		
		print "Write"
		read_defaults(bits)
		Din_s.write("0000 0000 0000 0000 0000 000")					#Din mux select
		write_en.write("0000 0000 0000 0000 1111 110")					#Row enable		
		if(inst[4]=="0"):								#precharge
		write_pre_0.write("0000 0000 0000 1111 1111 111")
		write_pre_1.write("1111 1111 1111 1111 1111 111")	
		write_col_en0.write("0000 0000 0000 1111 1111 111")			#column_block_select
		write_col_en1.write("0000 0000 0000 0000 0000 000")
	
		if(inst[4]=="1"):	
		write_pre_1.write("0000 0000 0000 1111 1111 111")
		write_pre_0.write("1111 1111 1111 1111 1111 111")	
		write_col_en1.write("0000 0000 0000 1111 1111 111")
		write_col_en0.write("0000 0000 0000 0000 0000 000")	
		
		for i in range(bits):								#Row address
			write_row_1.write("%s" % inst[5])
			write_row_0.write("%s" % inst[6])
			data1.write("%s" % inst[13])						#data
			data2.write("%s" % inst[14])
			data3.write("%s" % inst[15])
			data4.write("%s" % inst[16])		
	if(inst[0:4]=="0001"):									#Takes 30 cycles
		bits=17		
		print "Read"
		write_defaults(50)
		read_en.write("0000 0000 0000 0000 1000 0000 0000 00")						#Row enable		
		op_sel1.write("0000 0000 0001 1111 1111 1111 1111 11")
		if(inst[4]=="0"):								#precharge
		read_pre_0.write("0000 0000 0001 1111 1111 1111 1111 11")
		read_pre_1.write("1111 1111 1111 1111 1111 1111 1111 11")	
		block_en1.write(" 0000 0000 0001 1111 1111 1111 1111 11")				#column_block_select
		block_en2.write(" 0000 0000 0000 0000 0000 0000 0000 00")	
	
		if(inst[4]=="1"):	
		read_pre_1.write("0000 0000 0001 1111 1111 1111 1111 11")
		read_pre_0.write("1111 1111 1111 1111 1111 1111 1111 11")	
		block_en2.write(" 0000 0000 0001 1111 1111 1111 1111 11")
		block_en1.write(" 0000 0000 0000 0000 0000 0000 0000 00")	
		
		for i in range(11):								
			nor.write("0")	
			op_sel2.write("0")			
			r_add1.write("0")
			r_add0.write("0")
		
		for i in range(19):								#Row address
			nor.write("0")	
			op_sel2.write("0")			
			r_add1.write("%s" % inst[5])
			r_add0.write("%s" % inst[6])
			
		read_defaults(20)		

		def write_defaults(f):
			for i in range(f):							#Row address
				
					
			return

		def signal_gen(s,bit,f):
			
			for i in range(0,f):
				s = s + str(bit)	
			
			return s
			

		
#write_controls

#out
out_str = ""
out_str = signal_gen(out_str,1,43)
out_str = signal_gen(out_str,0,2)
out_str = signal_gen(out_str,1,43)
out.write(out_str)
#write_controls(out,out_str)

#enW
enW_str = ""
enW_str = signal_gen(enW_str,0,5)
enW_str = signal_gen(enW_str,1,9)
enW_str = signal_gen(enW_str,0,61)
enW.write(enW_str)
#write_controls(enW,enW_str)

#en1
en1_str = ""
en1_str = signal_gen(en1_str,0,1)
en1_str = signal_gen(en1_str,1,21)
en1.write(en1_str)
#write_controls(en1,en1_str)

#en
en_str = ""
en_str = signal_gen(en_str,0,7)
en_str = signal_gen(en_str,1,74)
en_str = signal_gen(en_str,0,70)
en.write(en_str)
#write_controls(en,en_str)

#enb
enb_str = ""
enb_str = signal_gen(enb_str,1,7)
enb_str = signal_gen(enb_str,0,74)
enb_str = signal_gen(enb_str,1,70)
enb.write(enb_str)
#write_controls(enb,enb_str)

#readP
readP_str = ""
readP_str = signal_gen(readP_str,0,13)
readP_str = signal_gen(readP_str,1,11)
readP_str = signal_gen(readP_str,0,70)
readP_str = signal_gen(readP_str,1,7)
readP_str = signal_gen(readP_str,0,50)
readP.write(readP_str)
#write_controls(readP,readP_str)

#RL
RL_str = ""
RL_str = signal_gen(RL_str,1,1)
RL_str = signal_gen(RL_str,0,21)
RL.write(RL_str)
#write_controls(RL,RL_str)

#writeP
writeP_str = ""
writeP_str = signal_gen(writeP_str,0,30)
writeP_str = signal_gen(writeP_str,1,11)
writeP_str = signal_gen(writeP_str,0,109)
writeP.write(writeP_str)
#write_controls(writeP,writeP_str)

#readP2
readP2_str = ""
readP2_str = signal_gen(readP2_str,0,13)
readP2_str = signal_gen(readP2_str,1,11)
readP2_str = signal_gen(readP2_str,0,126)
readP2.write(readP2_str)
#write_controls(readP2,readP2_str)

#S
S_str = ""
S_str = signal_gen(S_str,1,75)
S.write(S_str)
#write_controls(S,S_str)

#Sb
Sb_str = ""
Sb_str = signal_gen(Sb_str,1,75)
Sb.write(Sb_str)
#write_controls(Sb,Sb_str)


instruction_file.close()
write_row_0.close()
write_row_1.close()
write_pre_0.close()
write_pre_1.close()
write_col_en0.close()
write_col_en1.close()
write_en.close()
Din_s.close()
data1.close()
data2.close()
data3.close()
data4.close()

out.close()
enW.close()    
en1.close()    
en.close()	  
enb.close()	  
readP.close()	  
RL.close() 
writeP.close() 
readP2.close()	  
S.close()
Sb.close()

read_pre_0.close()
read_pre_1.close()
r_add0.close()    
r_add1.close()    
read_en.close()	  
op_sel1.close()	  
op_sel2.close()	  
block_en1.close() 
block_en2.close() 
nor.close()	  


\end{lstlisting} % filename in curly brackets

\clearemptydoublepage
%
% Add index
\printindex
%
\nocite{*}
\end{document} 